\cotefile{137}
\scanpage{73}
\uline{Ex 2} \quad Le ``complexe de De Rham \unsure{à|:} p. div''
est défini \deleted{comme} \unsure{canoniquement} (une \unsure{fois|fs}
connues les p.d.) comme \deleted{objet}
objet bigradué de $\deleted{\unsure{K}}\unsure{\mathcal R|\mathbb R|\mathcal Q}_{*}$, avec
structure d'algèbre bigraduée alternée
etc. Donc $F_C$ le transforme en
objet de la $=$ cat. On ne peut
affirmer encore (sans introduire
la ``structure p.d'' sur $F$) que
le complexe de De Rham associé à $F$
est alterné (seult antisym., i.e.,
\unsure{un signe|unité}). Le fait que
\deleted{\unsure{le} complexe de}
pour tt degré total donné ``$n$'', on trouve
une résolution de $h_*=\Psi_*^0$, flasque
(i.e. $\Sigma$-acyclique) sauf en degré
vert. maximum $n$ (où on trouve $\Psi_*^n
=\Lambda^n\Psi_*$) -- donc que c'est une résolution
$\Sigma$-acyclique tronquée, montre que
la \unsure{coho} de sa transformée par $F$
coïncide avec celle du $C^\bullet$ \deleted{\unsure{au fond}}\inserted{jusqu'en degré $n$}, et
\deleted{est nulle au dessus} -- mieux,
le complexe est homotope (avec
classe d'homotopie donnée) au
complexe tronqué \unsure{naïve|naïvem|naïf} de $C^\bullet$ --
en fait, on a un choix canonique.
(car $\overset{\inserted{\to\,\unsure{\deg\ \mathrm{tot}|\deg\leq n}}}{DR_*^{n\bullet}}\longrightarrow C^\bullet_*$ hom. \uline{unique}).
% archive sheet number: 72
